\documentclass[11pt]{article}
\newcommand{\shorttitle}{Squarefree products after deleting primes}
\usepackage[letterpaper,left=1.04in,right=1.04in,top=.92in,bottom=.9in,headheight=14pt]{geometry}
\usepackage[T1]{fontenc}
\usepackage{lmodern,amsmath,amssymb,amsthm,mathtools,booktabs,microtype,xcolor,fancyhdr,hyperref}
\definecolor{ink}{HTML}{173345}
\definecolor{accent}{HTML}{235E76}
\hypersetup{colorlinks=true,urlcolor=accent,linkcolor=accent,citecolor=accent}
\setlength{\emergencystretch}{1.5em}
\pagestyle{fancy}\fancyhf{}
\fancyhead[L]{\small\sffamily\color{accent}\shorttitle}
\fancyhead[R]{\small\sffamily\color{accent}26 November 2026}
\fancyfoot[C]{\small\thepage}
\renewcommand{\headrulewidth}{.3pt}
\newcommand{\papertitle}[2]{\thispagestyle{plain}\noindent{\small\sffamily\color{accent}#1 \hfill 26 November 2026}\par\vspace{.65em}{\raggedright\color{ink}\fontsize{24}{28}\selectfont\bfseries #2\par}\vspace{.65em}}
\newtheorem{theorem}{Theorem}
\newtheorem{lemma}[theorem]{Lemma}
\newtheorem{corollary}[theorem]{Corollary}
\newtheorem{proposition}[theorem]{Proposition}
\theoremstyle{remark}\newtheorem*{remark}{Remark}
\newcommand{\F}{\mathbb F}
\newcommand{\E}{\mathbb E}
\newcommand{\Prob}{\mathbb P}
\newcommand{\Span}{\operatorname{span}}
\hypersetup{pdftitle={Squarefree products after adversarial deletion of small primes}}
\begin{document}
\papertitle{Analytic number theory}{Squarefree products after adversarial deletion of small primes}
\noindent\textbf{Abstract.}
Take the primes up to $(\log q)^B$, with $B>2$, coprime to $q$. An adversary retains any proportion at least $1/2+\epsilon$. Products of exactly $k$ distinct retained primes are asymptotically equidistributed among the units modulo $q$ as soon as both $k$ and the number of primes omitted from the product exceed $C_\epsilon\log q$. The density threshold is sharp. More generally, retained density above $1/d$ dampens every character of order at least $d$. The analytic input is Dirichlet GRH via a known small-prime character estimate. The fixed-size selection and the uniformity in the deleted set follow from finite Fourier analysis.

\section{The arithmetic result}
Throughout this paper $\log$ denotes the natural logarithm. Let $B>2$, $q\ge3$, $x=(\log q)^B$, $G=(\mathbb Z/q\mathbb Z)^\times$, and
\[
 P=\{p\le x:p\text{ prime},\ p\nmid q\},\qquad N=|P|.
\]
Elements of $P$ are individual primes. Two different primes congruent modulo $q$ are distinct choices. For $A\subseteq P$, $M=|A|$, define
\[
 R_{A,k}(a)=\#\{S\subseteq A: |S|=k,\ \prod_{p\in S}p\equiv a\pmod q\}.
\]

\begin{theorem}[Uniform half-density threshold]\label{thm:main}
Fix $0<\epsilon<1/2$. For all sufficiently large $q$ (depending on $B,\epsilon$), every $A$ of size $M\ge(1/2+\epsilon)N$ and $1\le k<M$ satisfy, with $m=\min(k,M-k)$,
\[
 \max_{a\in G}\left|\frac{\varphi(q)R_{A,k}(a)}{\binom Mk}-1\right|
 \le 8\varphi(q)\sqrt{1+m}\,e^{-\epsilon m/2}.
\]
In particular, $m\ge(4/\epsilon)\log q$ gives $R_{A,k}(a)=(1+o(1))\binom Mk/\varphi(q)$ uniformly in $A,k,a$. For arbitrarily large $q$ there are retained sets of density $1/2+o(1)$ whose products stay in an index-two subgroup for every $k$.
\end{theorem}
Thus every reduced residue has many squarefree representatives with exactly the prescribed number of prime factors, even when the retained set is chosen with full knowledge of the modulus. Jao, Miller and Venkatesan study random walks generated by small primes~\cite{jmv}, and Pach studies fixed-cardinality subset sums drawn from a full finite abelian group~\cite{pach}. Here the allowed factors form an arbitrary dense subset of a sparse prime pool and each factor may appear once. The symmetric condition on $k$ comes from the bijection $S\mapsto A\setminus S$, which sends a product $a$ to $(\prod_{p\in A}p)/a$ in $G$.

\section{The analytic input}
For a nonprincipal Dirichlet character modulo $q$, the weighted prime estimate in Jao, Miller and Venkatesan~\cite[proof of Theorem 1.1, (2.6)]{jmv}, specialized to $\mathbb Q$, is
\[
 \sum_{p\le y}\chi(p)\log p\ll\sqrt y\,\log(2y)\log(2qy)\qquad(y\ge2).
\]
The implied constant is absolute. Their derivation uses GRH for the characters involved. The nonprincipal imprimitive case follows by passing to the inducing primitive character and removing primes dividing $q$, whose total logarithmic weight is at most $\log q$. Partial summation gives
\[
 \max_{\chi\ne1}\left|\sum_{p\in P}\chi(p)\right|\ll\sqrt x\log(2qx).
\]
Indeed the boundary term loses one $\log x$ and the remaining integral is bounded by a constant times $\log(2qx)\int_2^x dt/\sqrt t$. Also $N\sim x/\log x$ by the prime number theorem since at most $\log q/\log2$ primes divide $q$. Hence
\begin{equation}\label{eq:rho}
 \rho=\max_{\chi\ne1}\frac1N\left|\sum_{p\in P}\chi(p)\right|
 \ll_B\frac{\log\log q}{(\log q)^{B/2-1}}=o(1).
\end{equation}
It is the complex character sum in (2.6) that gives this bound for all phases. The proof needs only Dirichlet $L$-functions. The formal half-plane and open-strip theorems in \cite{lean} include these characters. The independent verification of the pinned edition is recorded in \cite{suth}.

\section{Selection without replacement}
\begin{lemma}\label{lem:fixed}
Let $z_1,\ldots,z_M$ be complex numbers of absolute value one and set $\Delta=1-M^{-1}|\sum_jz_j|$. For $1\le k<M$ and $m=\min(k,M-k)$,
\[
 \left|\binom Mk^{-1}\sum_{|S|=k}\prod_{j\in S}z_j\right|
 \le 8\sqrt{1+m}\exp\!\left(-\Delta\frac{k(M-k)}M\right).
\]
\end{lemma}
\begin{proof}
Set $u=k/(M-k)$. For every real $t$, the squared ratio
$|1+ue^{it}z_j|^2/(1+u)^2$ equals
$1-2u(1-\mathrm{Re}(e^{it}z_j))/(1+u)^2$. Using $1-v\le e^{-v}$ and multiplying yields
\[
 \prod_j |1+ue^{it}z_j|\le (1+u)^M
 \exp\!\left(-\frac{uM}{(1+u)^2}\Delta\right).
\]
Cauchy's coefficient bound gives the result up to the factor $(1+u)^M/(u^k\binom Mk)$. Its inverse is the probability of $k$ successes in $M$ Bernoulli trials with success probability $k/M$. The integer $k$ is a mode by the ratio of consecutive masses. The variance is at most $m$, so the interval of radius $2\sqrt{1+m}$ about $k$ carries at least $3/4$ of the mass by Chebyshev's inequality. It contains at most $5\sqrt{1+m}$ integers. The inverse modal probability is therefore at most $8\sqrt{1+m}$.
\end{proof}

\begin{proof}[Proof of Theorem~\ref{thm:main}]
For any $\chi\ne1$, deleting $N-M$ unit-complex summands gives $|\sum_{p\in A}\chi(p)|\le\rho N+N-M$. Consequently $\Delta\ge(2M-(1+\rho)N)/M\ge\epsilon$ for large $q$, by \eqref{eq:rho}. Apply Lemma~\ref{lem:fixed} to $z_p=\chi(p)$ and use $k(M-k)/M\ge m/2$. Character orthogonality expresses $\varphi(q)R_{A,k}(a)/\binom Mk-1$ as a sum of the nonprincipal normalized coefficients multiplied by $\overline{\chi(a)}$. Summing the bounds proves the inequality. For $m\ge(4/\epsilon)\log q$, the decreasing exponential times the square root is largest at the lower endpoint once $q$ is large. The relative error is $O_\epsilon(q^{-1}\sqrt{\log q})$.

For sharpness take $q$ through odd primes and retain exactly the primes with quadratic character $+1$. Their number is $(N+\sum_{p\in P}\chi(p))/2=(1/2+o(1))N$. Every product belongs to the index-two kernel.
\end{proof}

\section{Higher-order characters and the threshold \texorpdfstring{$1/d$}{1/d}}
\begin{theorem}\label{thm:higher}
Fix an integer $d\ge2$ and $0<\epsilon<1-1/d$. If $M\ge(1/d+\epsilon)N$, there exists $c_{d,\epsilon}>0$ such that for all sufficiently large $q$ and every character $\chi$ of order at least $d$,
\[
 \left|\binom Mk^{-1}\sum_{|S|=k}\chi\!\left(\prod_{p\in S}p\right)\right|
 \le 8\sqrt{1+m}\exp(-c_{d,\epsilon}m).
\]
The coefficient $1/d$ is sharp for a uniform assertion about every individual character of order at least $d$.
\end{theorem}
\begin{proof}
Choose $\theta$ so that the mean of $z_p=e^{-i\theta}\chi(p)$ on $A$ is real and nonnegative. On the unit circle put $F_d(z)=|1+z+\cdots+z^{d-1}|^2$.
Since $\chi^j$ is nonprincipal for $1\le j<d$, positivity and \eqref{eq:rho} give
\[
 \sum_{p\in A}F_d(z_p)\le\sum_{p\in P}F_d(z_p)
 \le N(d+d(d-1)\rho).
\]
The inequality $1-\mathrm{Re}(z^j)\le j^2(1-\mathrm{Re}\,z)$ follows by telescoping $1-z^j$ and squaring. Thus
\[
 F_d(z)\ge d^2-C_d(1-\mathrm{Re}\,z),\qquad
 C_d=2\sum_{j=1}^{d-1}(d-j)j^2=\frac{d^2(d^2-1)}6.
\]
These bounds force $\Delta\ge[d^2 M-N(d+d(d-1)\rho)]/(C_dM)\ge d^2\epsilon/(2C_d)$ for large $q$. The lemma gives $c_{d,\epsilon}=d^2\epsilon/(4C_d)$. For sharpness, take primes $q\equiv1\pmod d$, a character of exact order $d$, and retain its kernel in $P$. Averaging $1,\chi,\ldots,\chi^{d-1}$ gives retained density $1/d+o(1)$ and every product has character value one.
\end{proof}

For example, once the retained density exceeds $1/3$ by a fixed amount, every nonprincipal character except those of order two has an exponentially small coefficient. This identifies the characters that can still account for a substantial bias in a residue class. It makes no independence assumption about the chosen primes.

\clearpage
\section{A finite example and further questions}
For $q=101$, $x=251$ there are 53 eligible primes. Retain all 26 having quadratic character $+1$ and the nine smallest with character $-1$. Their $\binom{35}{14}=2{,}319{,}959{,}400$ products of fourteen distinct primes have the following exact counts over the 100 units.
\begin{center}
\begin{tabular}{lrr}\toprule
Retained set & Classes reached & Range of counts per unit\\\midrule
All 26 residues and nine nonresidues &100&23,197,877 to 23,201,277\\
Only the 26 residues &50&0 to 193,232\\\bottomrule
\end{tabular}
\end{center}
The example uses an independently specified cutoff to illustrate the deletion effect. Multiplying by each new prime and updating subset sizes in decreasing order computes the counts exactly. A direct enumeration in the accompanying script checks this recurrence on smaller inputs.

The diameter of the transition in $k$ is open in this argument. Does a sharper fixed-cardinality coefficient estimate give the best constant multiplying $\log q$ after a fixed deletion density? For density approaching $1/2$, can the distance from a quadratic fiber give a uniform replacement for $\epsilon$?

\begin{thebibliography}{9}
\bibitem{jmv} D. Jao, S. D. Miller, and R. Venkatesan, \href{https://doi.org/10.1016/j.jnt.2008.11.006}{Expander graphs based on GRH with an application to elliptic curve cryptography}, \emph{J. Number Theory} 129 (2009), 1491--1504. The proof of Theorem 1.1, (2.6), gives the prime estimate. \href{https://web.williams.edu/Mathematics/sjmiller/public_html/crypto/handouts/JaoMillerVenkatesan_ExpanderGraphsGRHEllCurveCrypto.pdf}{Author's PDF}.
\bibitem{pach} P. P. Pach, \href{https://arxiv.org/abs/2505.12496}{Equidistribution of subset sums}, arXiv:2505.12496 (2025).
\bibitem{lean} L. Sernau et al., \href{https://github.com/lsernau/grh-lean/blob/72fead102f663d5c7a407964a13dc0cf9d977207/grh/src/Main.lean}{GRH for Dirichlet $L$-functions, in Lean 4}, readable edition, commit \texttt{72fead10} (2026).
\bibitem{suth} A. V. Sutherland, \href{https://math.mit.edu/~drew/shortgrh/}{Verification files for the shorter Lean proof of RH and of GRH for Dirichlet $L$-functions} (26 November 2026).
\end{thebibliography}
\end{document}
